%-------------------------------------------------------------------------------
% SHARED STYLE
% Single source of truth for the visual design of both zpeng_cv.tex and
% zpeng_resume.tex. Edit here, and both documents stay in sync.
%-------------------------------------------------------------------------------

\usepackage{lastpage}
\usepackage{microtype}
\usepackage{needspace}

\RequirePackage[backend=biber,style=ieee,sorting=none]{biblatex}

%-------------------------------------------------------------------------------
% COLOR PALETTE
%
% One accent hue (matching the accent used on zpeng.me) plus a five-step
% neutral ramp. The ramp is what creates hierarchy: titles sit darkest,
% supporting metadata steps back, rules and the footer stay quiet.
%-------------------------------------------------------------------------------
\definecolor{awesome}{HTML}{1976D2}      % primary accent (zpeng.me blue)
\definecolor{accentdeep}{HTML}{14538D}   % deeper accent for the surname

\definecolor{darktext}{HTML}{22252A}     % entry / section titles
\definecolor{text}{HTML}{343A41}         % body copy
\definecolor{graytext}{HTML}{5C6570}     % roles, dates, secondary metadata
\definecolor{lighttext}{HTML}{8A929B}    % header location, footer
\colorlet{sectiondivider}{awesome!25!white}

% Keep the accent-colored first letters of each section heading.
\setbool{acvSectionColorHighlight}{true}

%-------------------------------------------------------------------------------
% TYPE SCALE
%
% Roughly a 1.15 modular scale off a 10pt body. The point of the scale is
% contrast: an entry title, its role, and its date must be visibly three
% different things. Weight and color do most of that work, size does the rest.
%-------------------------------------------------------------------------------

% Header
\renewcommand*{\headerfirstnamestyle}[1]{{\fontsize{26pt}{1em}\headerfontlight\color{graytext} #1}}
\renewcommand*{\headerlastnamestyle}[1]{{\fontsize{26pt}{1em}\headerfont\bfseries\color{accentdeep} #1}}
\renewcommand*{\headerpositionstyle}[1]{{\fontsize{9.5pt}{1.2em}\bodyfont\scshape\color{awesome} #1}}
\renewcommand*{\headeraddressstyle}[1]{{\fontsize{8.5pt}{1.2em}\headerfont\itshape\color{lighttext} #1}}
\renewcommand*{\headersocialstyle}[1]{{\fontsize{8.5pt}{1.2em}\headerfont\color{text} #1}}
\renewcommand*{\footerstyle}[1]{{\fontsize{8pt}{1em}\footerfont\scshape\color{lighttext} #1}}

% Sections
% \@sectioncolor is an internal of the class, so @ must be a letter here.
\makeatletter
\renewcommand*{\sectionstyle}[1]{{\fontsize{14pt}{1em}\bodyfont\bfseries\color{darktext}\@sectioncolor #1}}
\makeatother
\renewcommand*{\subsectionstyle}[1]{{\fontsize{10.5pt}{1em}\bodyfont\scshape\textcolor{graytext}{#1}}}
\renewcommand*{\paragraphstyle}{\fontsize{10pt}{1.38em}\bodyfontlight\upshape\color{text}}

% Entries
\renewcommand*{\entrytitlestyle}[1]{{\fontsize{11pt}{1.15em}\bodyfont\bfseries\color{darktext} #1}}
\renewcommand*{\entrypositionstyle}[1]{{\fontsize{9.5pt}{1.15em}\bodyfont\scshape\color{graytext} #1}}
\renewcommand*{\entrydatestyle}[1]{{\fontsize{9.5pt}{1.15em}\bodyfontlight\slshape\color{graytext} #1}}
\renewcommand*{\entrylocationstyle}[1]{{\fontsize{9.5pt}{1.15em}\bodyfontlight\slshape\color{awesome} #1}}
\renewcommand*{\descriptionstyle}[1]{{\fontsize{10pt}{1.35em}\bodyfontlight\upshape\color{text} #1}}

% Sub-entries
\renewcommand*{\subentrytitlestyle}[1]{{\fontsize{10pt}{1.15em}\bodyfont\mdseries\color{darktext} #1}}
\renewcommand*{\subentrypositionstyle}[1]{{\fontsize{9pt}{1.15em}\bodyfont\scshape\color{graytext} #1}}
\renewcommand*{\subentrydatestyle}[1]{{\fontsize{9pt}{1.15em}\bodyfontlight\slshape\color{graytext} #1}}
\renewcommand*{\subentrylocationstyle}[1]{{\fontsize{9pt}{1.15em}\bodyfontlight\slshape\color{awesome} #1}}
\renewcommand*{\subdescriptionstyle}[1]{{\fontsize{9.5pt}{1.3em}\bodyfontlight\upshape\color{text} #1}}

% Honors
\renewcommand*{\honortitlestyle}[1]{{\fontsize{9.5pt}{1.25em}\bodyfontlight\color{graytext} #1}}
\renewcommand*{\honorpositionstyle}[1]{{\fontsize{9.5pt}{1.25em}\bodyfont\bfseries\color{darktext} #1}}
\renewcommand*{\honordatestyle}[1]{{\fontsize{9.5pt}{1.25em}\bodyfontlight\color{graytext} #1}}
\renewcommand*{\honorlocationstyle}[1]{{\fontsize{9.5pt}{1.25em}\bodyfontlight\slshape\color{awesome} #1}}

% Skills
\renewcommand*{\skilltypestyle}[1]{{\fontsize{9.5pt}{1.3em}\bodyfont\bfseries\color{darktext} #1}}
\renewcommand*{\skillsetstyle}[1]{{\fontsize{9.5pt}{1.3em}\bodyfontlight\color{text} #1}}

%-------------------------------------------------------------------------------
% LAYOUT & RHYTHM
%-------------------------------------------------------------------------------
\geometry{left=2.2cm, top=1.8cm, right=2.2cm, bottom=1.9cm, footskip=.6cm}

% Long unbreakable strings in bibliography entries (paper titles, URLs) can
% otherwise push a line past the right margin.
\setlength{\emergencystretch}{2.5em}

\renewcommand{\acvHeaderAfterNameSkip}{1.2mm}
\renewcommand{\acvHeaderAfterPositionSkip}{1.4mm}
\renewcommand{\acvHeaderAfterAddressSkip}{0.6mm}
\renewcommand{\acvHeaderAfterSocialSkip}{4mm}
\renewcommand{\acvSectionTopSkip}{4.5mm}
\renewcommand{\acvSectionContentTopSkip}{2.8mm}

% Separator between the header's social links.
\renewcommand{\acvHeaderSocialSep}{\hspace*{0.6em}\textcolor{sectiondivider}{\textbar}\hspace*{0.6em}}

\makeatletter
% A lighter section rule, and a \needspace guard so a heading is never
% stranded alone at the foot of a page with its content overleaf.
\renewcommand{\cvsection}[1]{%
  \needspace{5\baselineskip}%
  \vspace{\acvSectionTopSkip}
  \sectionstyle{#1}
  \phantomsection
  \hspace{0.6em}\color{sectiondivider}\vhrulefill{1.1pt}%
}

% Subsections need a little air above them when they follow body text.
\renewcommand{\cvsubsection}[1]{%
  \needspace{4\baselineskip}%
  \vspace{\acvSectionContentTopSkip}
  \vspace{-2mm}
  \subsectionstyle{#1}
  \phantomsection
}
\makeatother

%-------------------------------------------------------------------------------
% LABELLED LISTS
%
% The stock cvskills table sizes its label column naturally and pins the value
% column to 0.9\textwidth, so any label wider than 0.1\textwidth pushes the
% value column off the right margin. Fixing the label column keeps the two
% together inside \textwidth no matter how long a label gets.
%-------------------------------------------------------------------------------
% It is also built on tabular*, which cannot break across pages, so a long
% list strands a half-empty page ahead of itself. longtable can break.
\usepackage{longtable}

\newlength{\cvskilllabelwidth}
\setlength{\cvskilllabelwidth}{3.9cm}
\newlength{\cvskillgap}
\setlength{\cvskillgap}{0.55cm}

\renewenvironment{cvskills}{%
  \vspace{\acvSectionContentTopSkip}
  \vspace{-2.0mm}
  \setlength\tabcolsep{0pt}
  \setlength{\extrarowheight}{0pt}
  \setlength{\LTpre}{0pt}
  \setlength{\LTpost}{0pt}
  \setlength{\LTleft}{0pt}
  \setlength{\LTright}{0pt}
  \begin{longtable}{@{} R{\cvskilllabelwidth} @{\hspace{\cvskillgap}} L{\textwidth - \cvskilllabelwidth - \cvskillgap} @{}}
}{%
  \end{longtable}
}

% Put a little air between consecutive rows of a skills/projects list.
\renewcommand*{\cvskill}[2]{%
  \skilltypestyle{#1} & \skillsetstyle{#2} \\[0.9mm]
}

% Tighten the gap after an entry that carries no description bullets
% (used by the education entries). Replaces ad-hoc \vspace{-1.2em} calls.
\newcommand{\cvnodescgap}{\vspace{-1.1em}}

% An accent-colored hyperlink, for the few places a URL is worth showing.
\newcommand{\cvurl}[2]{\href{#1}{\textcolor{awesome}{#2}}}
